% common/preamble.tex — paquetes compartidos por todos los talleres.
% Cargado desde cysec.cls. Motor: XeLaTeX.

\RequirePackage{etoolbox}

% --- Matemáticas: antes de unicode-math y de cleveref ---
% amssymb no se carga: choca con unicode-math. El ✓ de los mcq viene de pifont.
\RequirePackage{amsmath,mathtools}

% --- Fuentes (libertinus, igual que Tesis-MACC) ---
% Se cargan por nombre de archivo desde el árbol de TeX Live: las OTF de
% libertinus vienen con TeX Live pero normalmente NO están registradas en
% fontconfig, así que \setmainfont{Libertinus Serif} fallaría.
\RequirePackage{fontspec}
\RequirePackage{unicode-math}
\IfFontExistsTF{LibertinusSerif-Regular.otf}{%
  \setmainfont{LibertinusSerif}[
    Extension      = .otf,
    UprightFont    = *-Regular,
    ItalicFont     = *-Italic,
    BoldFont       = *-Bold,
    BoldItalicFont = *-BoldItalic]
  \setsansfont{LibertinusSans}[
    Extension   = .otf,
    UprightFont = *-Regular,
    ItalicFont  = *-Italic,
    BoldFont    = *-Bold]
  \setmonofont{LibertinusMono}[
    Extension   = .otf,
    UprightFont = *-Regular,
    Scale       = MatchLowercase]
  \setmathfont{LibertinusMath-Regular.otf}
}{%
  % Sin libertinus: se queda con Latin Modern, el default de XeLaTeX.
  \PackageWarning{cysec}{Fuentes Libertinus no encontradas; uso Latin Modern}%
}

% --- Idioma ---
% El idioma principal es el ÚLTIMO de la lista: babel lo reactiva en
% \begin{document}, así que un \selectlanguage en el preámbulo no bastaría
% (dejaba \today en inglés).
% shorthands=off desactiva los < > ~ activos de babel-spanish, que romperían
% cleveref; hace lo mismo que el modificador es-noshorthands, pero ese modificador
% no sobrevive al procesamiento de opciones con claves.
\ifdefstring{\cyseclang}{es}{%
  \RequirePackage[shorthands=off,english,spanish]{babel}%
}{%
  \RequirePackage[shorthands=off,spanish,english]{babel}%
}

% --- Color ---
\RequirePackage{xcolor}
\definecolor{accent}{HTML}{1F4E79}
\definecolor{oscrit}{HTML}{8B2500}
\definecolor{okgreen}{HTML}{1B7A3D}

% --- Listas, tablas, cajas ---
\RequirePackage{enumitem}
\RequirePackage{booktabs}
\RequirePackage{array}
\RequirePackage{longtable}
\RequirePackage{tabularx}
\RequirePackage[most]{tcolorbox}
\RequirePackage{graphicx}

% --- Código (comandos, payloads, sentencias) ---
% Estilo "cysec": mono de Libertinus (vía fontspec), keywords en accent y
% comentarios en gris, sin números de línea y con el mismo aire de caja
% sharp-corner que el resto de la clase.
\RequirePackage{listings}
\definecolor{codebg}{HTML}{F5F6F8}
\definecolor{codecom}{HTML}{5B6470}
\lstdefinestyle{cysec}{%
  basicstyle=\ttfamily\small,
  keywordstyle=\color{accent},   % sin \bfseries: Libertinus Mono no tiene negrita
  commentstyle=\color{codecom}\itshape,
  stringstyle=\color{oscrit},
  numberstyle=\tiny\color{codecom},
  numbers=none,
  backgroundcolor=\color{codebg},
  frame=single, framerule=0.4pt, rulecolor=\color{gray!45},
  framexleftmargin=4pt, framexrightmargin=4pt,
  xleftmargin=6pt, xrightmargin=6pt,
  breaklines=true, breakatwhitespace=true,
  showstringspaces=false,
  columns=fullflexible, keepspaces=true,
  aboveskip=0.8em, belowskip=0.8em,
}
\lstset{style=cysec}

% JavaScript no viene definido en listings (sí Java, HTML, SQL, bash...),
% así que se define aquí con el estándar ECMA-262 5th para los payloads web.
\lstdefinelanguage{JavaScript}{%
  keywords={break, case, catch, continue, debugger, default, delete, do,
    else, finally, for, function, if, in, instanceof, new, return, switch,
    this, throw, try, typeof, var, void, while, with, alert, document,
    window, true, false, null, undefined},
  sensitive=false,
  comment=[l]{//},
  morestring=[b]',
  morestring=[b]",
}

% --- Enlaces ---
\RequirePackage{xurl}          % quiebre de URLs largas en cualquier carácter
\RequirePackage{hyperref}
\hypersetup{
  colorlinks=true,
  linkcolor=accent, citecolor=accent, urlcolor=accent,
  pdfcreator={XeLaTeX + cysec.cls}
}
\RequirePackage[capitalise,noabbrev]{cleveref}
\ifdefstring{\cyseclang}{es}{%
  \crefname{figure}{Fig.}{Figs.}
  \crefname{table}{Tabla}{Tablas}
  \crefname{section}{Sección}{Secciones}
}{}

% --- Bibliografía compartida ---
\RequirePackage{csquotes}      % requerido por biblatex cuando hay babel
\RequirePackage[backend=biber,style=numeric,sorting=none,maxbibnames=99,
                giveninits=true,isbn=false,doi=false,url=true]{biblatex}
\addbibresource{bib/refs.bib}
\renewcommand{\bibfont}{\raggedright}   % URLs largas de las refs no sobresalen

\RequirePackage{microtype}
